% Компиляция: xelatex article30k_m2pro_ru.tex; bibtex article30k_m2pro_ru; xelatex ...; xelatex ...
\documentclass[11pt,a4paper]{article}

\usepackage[margin=2.35cm]{geometry}
\usepackage{fontspec}
% Arial is present on the target macOS system and provides Cyrillic bold/italic faces.
\setmainfont{Arial}
\usepackage[main=russian,english]{babel}
\usepackage{microtype}
\usepackage{amsmath}
\usepackage{booktabs}
\usepackage{tabularx}
\usepackage{array}
\usepackage{enumitem}
\usepackage{caption}
\usepackage{hyperref}
\usepackage{xurl}
\usepackage[numbers,sort&compress]{natbib}

\hypersetup{
  colorlinks=true,
  linkcolor=black,
  citecolor=black,
  urlcolor=blue,
  pdftitle={Мультитаск-классификация тональности и происхождения отзывов},
  pdfauthor={Автор указывается перед подачей}
}
\newcolumntype{Y}{>{\raggedright\arraybackslash}X}
\newcommand{\repo}[1]{\nolinkurl{#1}}
\newcommand{\en}[1]{\foreignlanguage{english}{#1}}
\setlength{\emergencystretch}{2em}

\title{Мультитаск-классификация тональности и происхождения отзывов с XLM-RoBERTa:\\
воспроизводимое исследование на частично размеченном корпусе из 30\,000 записей}
\author{\textbf{[ФИО автора]}\\
\textit{[образовательная программа, университет]}\\
\texttt{[author-email]}}
\date{Сентябрь 2026}

\begin{document}
\maketitle

\begin{abstract}
Онлайн-отзыв одновременно содержит оценочное суждение и сигнал о происхождении текста, однако эти свойства обычно моделируют раздельно. В работе исследуется, сохраняет ли совместное обучение общий языковой энкодер без потери качества на двух задачах: трёхклассовой классификации тональности и различении отзывов человека и отзывов, сгенерированных ИИ. Сформирован корпус из 30\,000 записей четырёх открытых наборов данных; разметка происхождения доступна только для 7\,000 гостиничных отзывов MAiDE-up. Сравниваются TF--IDF + логистическая регрессия, два одноцелевых и один мультитаск-вариант XLM-RoBERTa. Все варианты Transformer обучены на пяти начальных состояниях, а разбиение на обучение, валидацию и тест фиксировано и защищено от совпадений нормализованного текста между частями. На тесте мультитаск-модель достигла Macro-F1 $0.7709\pm0.0087$ для тональности и $0.9526\pm0.0094$ для происхождения; соответствующие показатели классического базиса составили $0.7161$ и $0.8998$. Разность с одноцелевыми Transformer-моделями статистически неубедительна: $-0.0013$ для тональности и $-0.0046$ для происхождения. Следовательно, в данном протоколе совместное обучение является компактной альтернативой двум моделям, но не демонстрирует доказанного положительного переноса относительно одноцелевого обучения. Результаты применимы к контролируемому внутридоменному протоколу и не являются доказательством универсальной детекции фальшивых отзывов.

\medskip
\noindent\textbf{Ключевые слова:} анализ тональности; мультитаск-обучение; XLM-RoBERTa; отзывы; тексты, сгенерированные ИИ; воспроизводимость.
\end{abstract}

\begin{center}
\textbf{Multi-task Classification of Review Sentiment and Provenance with XLM-RoBERTa: A Reproducible Study on a Partially Labelled Corpus of 30,000 Records}
\end{center}
\par\medskip

{\selectlanguage{english}
\renewcommand{\abstractname}{Abstract}
\begin{abstract}
Online reviews convey both evaluative sentiment and information about how a text was produced, but these properties are often modelled separately. This study tests whether a shared language encoder can address three-way sentiment classification and human-versus-AI review provenance without sacrificing predictive quality. A 30,000-record corpus was assembled from four open datasets; provenance labels are available only for 7,000 MAiDE-up hotel reviews. TF--IDF with logistic regression, two single-task XLM-RoBERTa systems, and a shared-encoder multi-task system were evaluated. Transformer variants were trained with five random seeds. The fixed train--validation--test partition was constructed by groups of normalized text, and the audit found no cross-partition text overlap. The multi-task model obtained test Macro-F1 of $0.7709\pm0.0087$ for sentiment and $0.9526\pm0.0094$ for provenance, compared with $0.7161$ and $0.8998$ for the classical baseline. Its differences from the corresponding single-task Transformer systems were inconclusive ($-0.0013$ and $-0.0046$ Macro-F1). Within this protocol, the shared model is therefore a compact alternative to two separate models, rather than evidence of positive transfer over matched single-task training. The provenance result is limited to the human-versus-AI hotel-review setting and must not be generalized to fake-review detection across domains.

\medskip
\noindent\textbf{Keywords:} sentiment analysis; multi-task learning; XLM-RoBERTa; online reviews; AI-generated text; reproducibility.
\end{abstract}
}

\section{Введение}

Тексты отзывов востребованы для мониторинга качества товаров и услуг, но их автоматический анализ сталкивается по меньшей мере с двумя различными вопросами. Первый вопрос --- какую оценку выражает автор. Второй --- относится ли текст к тому типу происхождения, для которого система была обучена. Появление генеративных моделей усиливает второй вопрос: правдоподобный отзыв может быть написан человеком или создан ИИ. Однако полярность текста и его происхождение не следует смешивать: отрицательная оценка не делает отзыв ``подозрительным'', а высокий классификационный балл не устанавливает факт обмана.

В настоящей работе происхождение определено операционально как бинарная метка набора MAiDE-up: человеческий отзыв против ИИ-сгенерированного гостиничного отзыва. MAiDE-up содержит реальные и ИИ-сгенерированные гостиничные отзывы в десяти языках; его авторы специально рассматривают зависимость качества распознавания от языка, локации и тональности~\cite{ignat2025maide}. Поэтому более широкие формулировки вроде «детекция фейковых отзывов» здесь намеренно не используются.

Мультитаск-обучение предполагает, что связанные задачи используют общее представление и тем самым могут служить друг для друга индуктивным смещением~\cite{caruana1997multitask}. Для многоязычного корпуса естественным кандидатом на общий энкодер является XLM-RoBERTa (XLM-R), предобученная модель на 100 языках~\cite{conneau2020xlmr}. Но сама по себе общая архитектура не гарантирует переноса: добавочная задача может не помочь, а сравнение с одноцелевым вариантом необходимо для отделения эффекта общего энкодера от эффекта более сильной модели.

Цель исследования --- воспроизводимо сравнить три семейства классификаторов на едином корпусе и аккуратно ограничить выводы его дизайном. Проверяются следующие вопросы:
\begin{enumerate}[label=\textbf{RQ\arabic*.}, leftmargin=1.5cm]
  \item Улучшает ли мультитаск-XLM-R Macro-F1 относительно TF--IDF + логистической регрессии на каждой задаче в фиксированном тестовом протоколе?
  \item Отличается ли мультитаск-XLM-R от одноцелевого XLM-R, обученного для той же задачи?
  \item Какие свойства данных и эксперимента ограничивают перенос результатов за пределы данного протокола?
\end{enumerate}

Вклад работы состоит не в заявлении нового универсального детектора, а в документированном эксперименте: приведены происхождение и состав данных, правила разбиения, конфигурация обучения, многосидовые результаты, парные статистические сравнения и ограничения валидности.

\section{Связанные работы}

Идея совместно обучать несколько задач с общими параметрами восходит к постановке индуктивного переноса через общую репрезентацию~\cite{caruana1997multitask}. В прикладной постановке это означает, что один энкодер получает градиенты от нескольких голов классификации. Польза такого решения является эмпирическим, а не аксиоматическим вопросом: эффективность зависит от родства задач, объёма разметки, баланса классов и процедуры обучения.

XLM-R предназначена для переноса между языками и была предобучена как многоязычная masked-language model на данных ста языков~\cite{conneau2020xlmr}. Это мотивирует её использование в корпусе, где 23\,734 из 30\,000 записей русскоязычны, а остальные представлены девятью языками. Такая мотивация не равна доказательству межъязыковой устойчивости: в текущем исследовании языковые срезы не являются независимым внешним тестом.

Для задачи происхождения важен характер исходной разметки. MAiDE-up создан именно для различения настоящих и ИИ-сгенерированных гостиничных отзывов и содержит по 10\,000 примеров каждого типа в десяти языках~\cite{ignat2025maide}. В данной работе из него использована квотированная подвыборка 7\,000 текстов. Таким образом, результат по происхождению относится к типу синтетичности и домену, заданным этим ресурсом, а не к любым человеческим обманным практикам.

Наконец, сравнение моделей на одном тесте требует не только среднего значения метрики. Рекомендации по статистическому тестированию в NLP подчёркивают зависимость выбора теста от экспериментального дизайна и фиксированного набора примеров~\cite{dror2018hitchhiker}. Поэтому далее отдельно сообщаются вариация по начальным состояниям, доверительные интервалы разностей по тестовым примерам и приближённый randomization test.

\section{Данные и метод}

\subsection{Корпус и метки}

Входной файл \repo{data/processed/joint\_reviews.article30k.jsonl} содержит 30\,000 нормализованных записей. Его SHA-256 равен \repo{a296cf7bb8cb0dbb3eebd8a0cc9c478422e9f5959836bc4d20915f9f81c9c5fc}. Источники и их роли сведены в табл.~\ref{tab:data}. RuReviews, карточка Perekrestok и карточка Wildberries использовались как открытые источники русскоязычных потребительских отзывов~\cite{rureviews,perekrestok,wildberries}; MAiDE-up предоставил двоичную разметку происхождения~\cite{ignat2025maide}.

\begin{table}[ht]
\centering
\caption{Состав корпуса и происхождение разметки.}
\label{tab:data}
\small
\begin{tabularx}{\linewidth}{@{}l r Y Y@{}}
\toprule
Источник & $n$ & Тональность & Происхождение \\
\midrule
MAiDE-up & 7\,000 & Метка исходного набора; классы negative/positive & 3\,500 human (authentic) и 3\,500 AI-generated (fake) \\
Perekrestok & 10\,000 & Эвристика рейтинга: 1--2 negative, 3 neutral, 4--5 positive & Нет метки \\
RuReviews & 3\,000 & Исходная трёхклассовая метка набора & Нет метки \\
Wildberries & 10\,000 & Эвристика рейтинга: 1--2 negative, 3 neutral, 4--5 positive & Нет метки \\
\bottomrule
\end{tabularx}
\end{table}

Все 30\,000 записей имеют метку тональности: 23\,734 текста на русском языке и 6\,266 на турецком, французском, корейском, английском, китайском, испанском, итальянском, немецком и румынском языках. Доменный состав: 23\,000 e-commerce и 7\,000 hospitality. Метка происхождения присутствует только у 7\,000 записей MAiDE-up. Поэтому мультитаск-потери вычисляются только по доступным меткам: e-commerce записи участвуют в обучении тональности, но не добавляют наблюдений к функции потерь происхождения.

Две части тональной разметки получены из рейтинга, а не из независимой ручной аннотации. Это сознательный компромисс между объёмом материала и точностью целевой переменной. Он должен учитываться при интерпретации: измеряется согласие с трёхклассовым преобразованием рейтинга, а не чистая лингвистическая полярность во всех возможных значениях этого термина.

\subsection{Разбиение и контроль утечки}

Корпус разделён с фиксированным seed 42 на 23\,975 обучающих, 3\,013 валидационных и 3\,012 тестовых записей. Перед разделением тексты приводились к casefold-представлению с нормализацией пробелов и объединялись в группы по нормализованному тексту. Группы стратифицировались по источнику и паре доступных меток. Следовательно, одинаковый нормализованный текст не может попасть в разные части.

Аудит действительно нашёл 521 точную и 564 нормализованные дублирующиеся строки во всём объединённом корпусе, но нулевые пересечения между train/validation/test одновременно по идентификатору, точному тексту и нормализованному тексту. Наличие дублей остаётся характеристикой качества исходного материала, однако не создаёт наблюдаемой текстовой утечки через данное разбиение. В тестовой части распределение тональности составляет 710 negative, 204 neutral и 2\,098 positive; для происхождения оно сбалансировано (350 authentic и 350 fake). По этой причине основной метрикой выбрана Macro-F1, а accuracy приводится дополнительно.

\subsection{Модели}

Базовый классификатор состоит из TF--IDF признаков уни- и биграмм (не более 50\,000 признаков) и отдельной \en{LogisticRegression} с балансировкой классов для каждой доступной задачи. Он служит сильной неглубокой точкой отсчёта на тех же обучающих текстах.

Одноцелевые и мультитаск-модели используют \repo{FacebookAI/xlm-roberta-base}. В мультитаск-варианте последний скрытый вектор первого токена (CLS) подаётся в две независимые MLP-головы: одну для трёх классов тональности и одну для двух классов происхождения. Каждая голова имеет схему $h\rightarrow h\rightarrow k$, нелинейность GELU и dropout 0.1. Потеря равна сумме взвешенных кросс-энтропий доступных задач; веса обеих задач заданы равными 0.5. Одноцелевые варианты используют тот же энкодер и MLP-голову, но оптимизируют только свою метку.

Конфигурация ориентирована на доступную локальную среду Apple Silicon: максимальная длина 192 токена, batch size 2, накопление градиента на четырёх шагах, gradient checkpointing, AdamW с learning rate $2\cdot10^{-5}$, weight decay 0.01, линейный warmup 10\%, максимум пять эпох и early stopping с patience 2. В мультитаск-обучении используется source-balanced sampler и взвешенная кросс-энтропия классов; в одноцелевых запусках source-balanced sampler не применяется. Это означает, что сравнение ``multitask против single-task'' оценивает два настроенных учебных контура, а не изолированно один архитектурный переключатель.

Обучение выполнено на Mac mini M2 Pro с 16~GB unified memory через backend MPS. Записанные среда и конфигурация: Python 3.12.12, macOS 26.2 arm64, устройство MPS, пять train seeds $\{11,21,42,84,126\}$. Выбор чекпойнта делался по валидационному Macro-F1, тогда как тест использовался для итоговой оценки.

\subsection{Протокол сравнения}

Для каждого seed обучались мультитаск-модель и два одноцелевых Transformer-варианта. Классический базис воспроизводимо обучается детерминированно при фиксированной конфигурации, поэтому его стандартное отклонение по повторённым записям равно нулю; это не следует трактовать как общую меру стабильности метода. Для каждого Transformer сообщаются среднее и стандартное отклонение по пяти seed, а также t-интервал 95\% для средних метрик.

Парные сравнения выполнялись на одинаковых тестовых примерах: bootstrap из 2\,000 выборок для доверительного интервала разности, усреднённой по seed, и approximate randomization test из 2\,000 перестановок. Восемь проверок (две задачи, две метрики, два сопоставления) должны читаться как семейство связанных тестов; поэтому главным основанием вывода служат направление разности, интервал и консервативная формулировка, а не одно p-значение.

\section{Результаты}

\subsection{Качество на тестовой части}

Таблица~\ref{tab:mainresults} содержит итоговые показатели. Мультитаск-XLM-R превосходит TF--IDF + логистическую регрессию по Macro-F1 в обеих задачах. Для тональности прирост составляет 0.0548, для происхождения --- 0.0528. Одноцелевые Transformer-модели имеют немного более высокие средние значения, однако по одному этому наблюдению нельзя заключать о преимуществе: разброс по seed сопоставим с разностью средних.

\begin{table}[ht]
\centering
\caption{Тестовые метрики; для Transformer указано среднее $\pm$ стандартное отклонение по пяти seed.}
\label{tab:mainresults}
\small
\begin{tabularx}{\linewidth}{@{}Y Y r r@{}}
\toprule
Семейство & Задача & Macro-F1 & Accuracy \\
\midrule
TF--IDF + LR & Тональность & 0.7161 & 0.8320 \\
Single-task XLM-R & Тональность & $0.7721\pm0.0087$ & $0.8958\pm0.0038$ \\
Multitask XLM-R & Тональность & $0.7709\pm0.0087$ & $0.8950\pm0.0047$ \\
\addlinespace
TF--IDF + LR & Происхождение & 0.8998 & 0.9000 \\
Single-task XLM-R & Происхождение & $0.9571\pm0.0088$ & $0.9571\pm0.0087$ \\
Multitask XLM-R & Происхождение & $0.9526\pm0.0094$ & $0.9526\pm0.0094$ \\
\bottomrule
\end{tabularx}
\end{table}

Для тональности t-интервал Macro-F1 мультитаск-модели равен $[0.7600;0.7817]$, а одноцелевой --- $[0.7613;0.7830]$. Для происхождения интервалы соответственно равны $[0.9409;0.9642]$ и $[0.9463;0.9680]$. Эти интервалы описывают вариацию и неопределённость по пяти начальным состояниям при одном и том же разбиении, а не неопределённость выбора новых источников или новых наборов текстов.

\subsection{Парные сравнения}

Результаты основных сравнений Macro-F1 представлены в табл.~\ref{tab:comparisons}. При сравнении с классическим базисом 95\% bootstrap-интервалы положительны, а приближённые p-значения равны 0.0005 (минимально разрешимое значение при 2\,000 перестановках). При сравнении с одноцелевым Transformer оба интервала пересекают нуль, а p-значения не дают основания отвергнуть отсутствие различия. В частности, мультитаск-вариант не доказал положительный перенос, поскольку его среднее значение немного ниже в обеих задачах.

\begin{table}[ht]
\centering
\caption{Парные сравнения мультитаск-модели по Macro-F1. $\Delta=A-B$; интервал построен bootstrap по общему тесту.}
\label{tab:comparisons}
\small
\begin{tabularx}{\linewidth}{@{}Y Y r Y r@{}}
\toprule
Задача & Сравнение ($A$ против $B$) & $\Delta$ & 95\% CI & $p_{AR}$ \\
\midrule
Тональность & multitask -- single-task & $-0.0013$ & $[-0.0105;0.0081]$ & 0.8016 \\
Тональность & multitask -- TF--IDF+LR & $+0.0548$ & $[0.0337;0.0732]$ & 0.0005 \\
Происхождение & multitask -- single-task & $-0.0046$ & $[-0.0109;0.0020]$ & 0.1829 \\
Происхождение & multitask -- TF--IDF+LR & $+0.0528$ & $[0.0305;0.0749]$ & 0.0005 \\
\bottomrule
\end{tabularx}
\end{table}

Таким образом, ответ на RQ1 положителен в рамках фиксированного тестового протокола: мультитаск-XLM-R показывает более высокий Macro-F1, чем использованный классический базис. Ответ на RQ2 иной: данных недостаточно, чтобы заявить различие мультитаск- и одноцелевого XLM-R. Практическая интерпретация ограничена следующей: общий энкодер позволяет обслуживать две головы одной моделью без заметной потери среднего качества по отношению к двум одноцелевым Transformer-запускам, но не доказал превосходства над ними.

\section{Обсуждение и ограничения}

Полученная картина асимметрична не по задачам, а по уровню сравнения. Transformer-варианты устойчиво опережают выбранный TF--IDF-базис, тогда как разность между общей и отдельной архитектурами мала относительно оценки неопределённости. Этот результат совместим с несколькими объяснениями: задачи могут использовать частично общие признаки, но объём и тип меток происхождения недостаточны, чтобы сделать выигрыш от совместного обучения измеримым. Эксперимент не позволяет выбирать между этими объяснениями, поэтому причинное утверждение о механизме переноса не делается.

У исследования есть следующие существенные ограничения.
\begin{enumerate}[leftmargin=0.65cm]
  \item \textbf{Ограниченная валидность происхождения.} Все 700 тестовых меток происхождения происходят из MAiDE-up и относятся к гостиничному домену, где противопоставлены человеческие и ИИ-сгенерированные отзывы. Это не внешний тест для e-commerce и не тест человеческих обманных отзывов.
  \item \textbf{Неоднородная семантика тональных меток.} Для Perekrestok и Wildberries тональность получена из рейтинга. Следовательно, доменные различия могут отражать не только язык и стиль текста, но и процедуру формирования меток.
  \item \textbf{Один фиксированный split.} Пять seed измеряют чувствительность обучения к инициализации, но не заменяют повторные разбиения, leave-one-source-out или внешний набор. Следующая сильная проверка должна удерживать целый источник, домен или язык вне обучения.
  \item \textbf{Дубликаты и отбор.} Групповое разбиение исключило наблюдаемую межсплитовую утечку, однако 564 нормализованных дублирующихся строки остаются свойством объединённого корпуса. Кроме того, источники включались в заданных квотах, поэтому агрегатная метрика не равна оценке случайной популяции отзывов.
  \item \textbf{Неидеальная абляция.} В мультитаск-режиме включён source-balanced sampler, в одноцелевом --- нет. Поэтому сравнение не является чистой абляцией только общей головы; для такого вывода требуется повторить опыты при полностью одинаковом sampling schedule.
  \item \textbf{Нет калибровки и эксплуатационной валидации.} Вероятности softmax не калибровались, порог человеческой проверки не выбирался, а модельные артефакты не были автоматически назначены production-моделью приложения. Следовательно, показатели не задают готовую политику модерации.
\end{enumerate}

Эти ограничения определяют корректную дорожную карту. Во-первых, необходимо собрать независимые метки происхождения как минимум для ещё одного домена и для человеческой deception-разметки. Во-вторых, следует повторить сравнение на source-held-out и language-held-out протоколах с одинаковым sampler для multitask и single-task. В-третьих, для прикладного решения нужны калибровка, human-in-the-loop интерфейс и анализ ошибок, а не только усреднённый F1.

\section{Заключение}

На частично размеченном смешанном корпусе из 30\,000 отзывов реализован и проверен воспроизводимый эксперимент с общим XLM-RoBERTa-энкодером для тональности и происхождения текста. Мультитаск-модель повысила тестовый Macro-F1 относительно TF--IDF + логистической регрессии с 0.7161 до 0.7709 для тональности и с 0.8998 до 0.9526 для происхождения. При этом одноцелевые Transformer-модели показали практически сопоставимые, немного более высокие средние значения, а парные интервалы не подтверждают различия.

Следовательно, результат поддерживает использование общей модели как компактного исследовательского артефакта, но не поддерживает тезис о её доказанном преимуществе перед одноцелевым Transformer или о всеобщей детекции фальшивых отзывов. Наиболее важное продолжение --- внешний, межисточниковый и междоменный тест с независимой разметкой.

\section*{Доступность данных и воспроизводимость}

Скрипты, конфигурация и агрегированные отчёты доступны в локальном репозитории проекта: \repo{configs/model.article30k.mac\_m2\_16gb.yaml}, \repo{reports/multiseed/article30k-m2pro-retry1/}, \repo{reports/multiseed/article30k-m2pro-retry1\_statistics/} и \repo{reports/audit/joint\_reviews.article30k.audit.json}. Контрольная сумма обработанного корпуса приведена в разделе~3.1. Повторное распространение исходных записей определяется условиями исходных наборов; в частности, Wildberries-компонент имеет лицензию CC BY-NC-SA 4.0. Для строгого воспроизведения следует использовать зафиксированный конфигурационный снимок и опубликованные карточки наборов данных, а не подменять их новыми версиями.

\section*{Этическая декларация}

В работе не проводился набор участников и не собирались новые персональные данные. Анализировались текстовые записи, полученные из открытых исследовательских источников; в рукописи не воспроизводятся сырые отзывы, имена авторов или прямые идентификаторы пользователей. Предсказание происхождения следует использовать как исследовательский сигнал, а не как самостоятельное доказательство недобросовестности автора.

\section*{Вклад авторов по CRediT}

\textbf{Перед подачей необходимо заменить этот шаблон реальными именами и ролями.} Для одноавторской версии ожидаемая формулировка: \textit{[Автор]: концептуализация, методология, программное обеспечение, подготовка данных, формальный анализ, валидация, визуализация, написание первоначального текста, редактирование.}

\section*{Финансирование и конфликт интересов}

Информация о финансировании и конфликте интересов автором не предоставлена. Перед подачей необходимо явно заменить этот текст одной из деклараций: «Финансирование отсутствует» или указать грант; «Авторы заявляют об отсутствии конфликта интересов» либо описать существующий конфликт.

\section*{Раскрытие использования ИИ}

Генеративный ИИ использовался как вспомогательный инструмент при подготовке структуры и языковой версии рукописи. Численные результаты, конфигурация эксперимента, утверждения о данных и ссылки были сверены с артефактами репозитория и первичными источниками. Автор несёт полную ответственность за окончательную проверку рукописи, трактовку результатов и соблюдение требований площадки подачи.

\bibliographystyle{unsrtnat}
\bibliography{article30k_m2pro_ru}

\end{document}
